\section{Ablation: what each ingredient contributes}\label{sec:ablation}

The argument has four removable ingredients: the restoration lemma, each of the two computed transmitted sets, the transfer relation
over the charge classes, and the antichain form of the closure. Table~\ref{tab:ablation} states what remains when each is removed; the subsections
below read the table.

\begin{table}[t]\centering\footnotesize
\begin{tabular}{p{4.2cm}p{8.2cm}}
\toprule
ingredient removed & what remains \\
\midrule
restoration Lemma~\ref{lem:restore} & the theorem still holds: the refutations of distance 0 for $A$ (7.7 s) and of distance 1 for $B$ (2.8 s) replace it \\
computation on $B$ & $D(A) \subseteq \{1\}$ by Lemma~\ref{lem:restore}(a) and two refutations (2.7 s, 2.1 s), and $1 \notin D(B)$ by Lemma~\ref{lem:restore}(b): disjoint \\
computation on $A$ & $D(B) \subseteq \{0\}$ by Lemma~\ref{lem:restore}(b) and two refutations (2.2 s, 2.1 s), and $0 \notin D(A)$ by Lemma~\ref{lem:restore}(a): disjoint \\
transfer relation, classes only & the conducted sets of the two sides of the 6-edge cut of $\Gfty$ share $\{\Oz,\OEc,\ODtwo\}$ in all ten splits, so no argument by classes separates them \\
antichain form, exact closure instead & $F_2$: 19{,}005 elements instead of 276 (a factor of about 69); $F_3$: 116{,}463 instead of 2{,}454 (about 47); $F_4$: stopped at the cap \\
\bottomrule
\end{tabular}
\caption{Ablation of the argument: the ingredient removed and what remains, from the transmitted sets of Table~\ref{tab:sets}, the control of Table~\ref{tab:control} and the closure sizes of Table~\ref{tab:closures}. Times are solver times of computation EXP-012.}\label{tab:ablation}
\end{table}

\subsection{The two transmitted sets and the restoration lemma}\label{sec:ablation-sets}

The disjointness $D(A) \cap D(B) = \emptyset$ is proved three times over. The restoration lemma excludes distance 0 for $A$ and distance 1 for $B$
(Lemma~\ref{lem:restore}); computation excludes the same two distances by refutations of 7.7 and 2.8 seconds; and computation
excludes the two remaining distances for each block. Any one of the three subsets of the argument that removes one ingredient still proves
Theorem~\ref{thm:main}, as Table~\ref{tab:ablation} shows. This redundancy is a check on the encoder: the computation could have
returned a map for distance 0 on $A$ or for distance 1 on $B$, which would have contradicted the restoration lemma and exposed an error in the lemma or
in the formula. It did not. The three refutations for $A$ take 12.5 seconds in total and the three for $B$ 7.2 seconds; the smallest set of computations that proves
the theorem together with the restoration lemma is the pair of refutations of one block, 4.9 seconds for $A$ (distances 2 and 3) or 4.4 seconds for $B$, so each pair is by itself a complete second proof of the disjointness.

\subsection{Classes without the transfer relation}\label{sec:ablation-control}

The graph $\Gfty$ has a 6-edge cut with sides of 18 and 34 vertices (cut edges 8, 26, 39, 40, 48, 59). For each of the ten splits of the cut into two triples,
the two sides are 6-poles with Petersen colorings whose conducted sets share $\Oz$, $\OEc$ and $\ODtwo$, so no argument by classes separates them
(Table~\ref{tab:control}). Their ring of length 2 is $\Gfty$, which has no Petersen coloring, and the product of their sector matrices has a zero diagonal in all six
sectors in every split. The classes therefore do not detect the non-colorability of $\Gfty$, while the transfer relation does: this is the evidence that the
semigroup test is not vacuous, in the sense that it returns a zero diagonal on a gluing of two colorable 6-poles that is known not to be colorable.

\begin{table}[t]\centering\footnotesize
\setlength{\tabcolsep}{4pt}
\begin{tabular}{rlll}
\toprule
split & cut edges on the first connector & conducted by side $X$ & conducted by side $V - X$ \\
\midrule
1 & 8, 26, 39 & $\{\Oz,\OEc,\ODthree,\ODtwo,\OTtwo\}$ & $\{\Oz,\OEc,\ODthree,\ODtwo\}$ \\
2 & 8, 26, 40 & $\{\Oz,\OEc,\ODthree,\ODtwo,\OTtwo\}$ & $\{\Oz,\OEc,\ODthree,\ODtwo\}$ \\
3 & 8, 26, 48 & $\{\Oz,\OEc,\ODtwo\}$ & $\{\Oz,\OEc,\ODthree,\ODtwo\}$ \\
4 & 8, 26, 59 & $\{\Oz,\OEc,\ODtwo\}$ & $\{\Oz,\OEc,\ODtwo\}$ \\
5 & 8, 39, 40 & $\{\Oz,\OEc,\ODtwo\}$ & $\{\Oz,\OEc,\ODthree,\ODtwo\}$ \\
6 & 8, 39, 48 & $\{\Oz,\OEc,\ODthree,\ODtwo,\OTone,\OTtwo\}$ & $\{\Oz,\OEc,\ODthree,\ODtwo,\OTtwo\}$ \\
7 & 8, 39, 59 & $\{\Oz,\OEc,\ODthree,\ODtwo\}$ & $\{\Oz,\OEc,\ODthree,\ODtwo,\OTtwo\}$ \\
8 & 8, 40, 48 & $\{\Oz,\OEc,\ODthree,\ODtwo,\OTone,\OTtwo\}$ & $\{\Oz,\OEc,\ODthree,\ODtwo,\OTtwo\}$ \\
9 & 8, 40, 59 & $\{\Oz,\OEc,\ODthree,\ODtwo\}$ & $\{\Oz,\OEc,\ODthree,\ODtwo,\OTtwo\}$ \\
10 & 8, 48, 59 & $\{\Oz,\OEc,\ODthree,\ODtwo\}$ & $\{\Oz,\OEc,\ODthree,\ODtwo,\OTtwo\}$ \\
\bottomrule
\end{tabular}
\caption{The ten splits of the 6-edge cut of $\Gfty$ (cut edges 8, 26, 39, 40, 48, 59; sides of 18 and 34 vertices) into two triples, with the orbits conducted by each side (computation EXP-014, \texttt{g52-control.json}). The sides always share $\Oz$, $\OEc$ and $\ODtwo$. For every split the product of the two sector matrix tuples, which is the ring of length 2 and equals $\Gfty$, has a zero diagonal in all six sectors.}\label{tab:control}
\end{table}


\subsection{What the two comparisons with the graphs test}\label{sec:ablation-direction}

A computed sector matrix contains only entries that are checked maps, so by Lemma~\ref{lem:soundness} a nonzero diagonal of a product is
always correct. The comparison with the solver can therefore expose a missing entry (a zero trace on a ring for which the solver finds a coloring) or an error in the
construction of the ring graph (a nonzero trace on a ring that the solver finds not colorable). The \nCrossRings{} agreements of Section~\ref{sec:crosscheck} are all of the kind ``nonzero trace and
coloring found'', since every ring was colorable; they show that the computed matrices miss nothing that these rings need, and they cannot test the converse
direction, a zero trace on a ring that the solver finds not colorable, because no such ring occurred. The control of Table~\ref{tab:control}
supplies ten cases of that converse direction, all on the one graph $\Gfty$. The two comparisons together cover both directions only on those cases. We therefore state
which claims need which direction.

\begin{itemize}
\item Claims that rest on soundness only: every ring of claws, of superedges and of both is colorable (Theorem~\ref{thm:ringsY}); every ring of the
$\PG$, $J_5$, $J_7$ and dodecahedron blocks that has a nonzero trace is colorable.
\item Claims that rest on the completeness of the computed matrices: the sector of each flower snark ring (Corollary~\ref{cor:flower}), which is supported for the claw by the exhaustive
enumeration and the second encoding, and for $k = 5, 7, 9$ by the explicit rings; the zero diagonals of the control, which are consistent with the known
non-colorability of $\Gfty$ but would also appear if matrices were incomplete; and any statement that a sector is not conducted, which for the 6-poles of Table~\ref{tab:splits} is
decided by the formulas of computation EXP-013 and not by the matrices.
\item Claims that no computation settles: the existence of a zero-trace element in the closures of $F_4$ and $F_5$ (Table~\ref{tab:closures}).
\end{itemize}

\subsection{Exact closure against antichain}\label{sec:ablation-closure}

The antichain form changes the cost differently by family. For $F_2$ the certificate has 276 elements against 19{,}005 and is computed in 5.0 seconds against 27.1; for
$F_3$ the certificate is 47 times smaller (2{,}454 against 116{,}463) but takes 349.4 seconds against 344.6, so the containment tests offset the saving, whose cost is not measured separately; for $F_1$ both take under a second. For $F_4$ the exact closure stopped at the cap after 93.4 seconds and the antichain closure was stopped at the one-hour limit; the 336 generators give $336^2 = 112{,}896$ candidate products of two generators, so at least $200{,}070 - 112{,}896 = 87{,}174$ of the stored elements are words of length 3, which is how far the exact closure advanced before the cap. So on the
larger families the exact closure was stopped by the cap and the antichain closure by time. A faster containment test, for example by indexing the kept elements,
is a possible route to the two undecided families \claim{C}, and is not part of this work. All these comparisons are made on the same generators, the same cap of 200{,}000 elements.
